\documentclass[12pt, a4paper]{article}

\usepackage[utf8]{inputenc}
\usepackage[T1]{fontenc}
\usepackage[brazil]{babel}
\usepackage{amsmath, amssymb}
\usepackage[margin=2.5cm]{geometry}

\title{Formulação do Equivalente Determinístico Multiestágio}
\date{}

\begin{document}

\maketitle

Para a implementação do Equivalente Determinístico (DEQ) em formato de árvore de cenários, as variáveis e restrições são indexadas por nós da árvore, preservando a dependência intertemporal e a não-antecipatividade.

\section{Conjuntos e Índices}
\begin{itemize}
    \item $\mathcal{N}$: Conjunto de todos os nós da árvore de cenários.
    \item $\mathcal{L} \subset \mathcal{N}$: Subconjunto dos nós folha (último estágio da árvore).
    \item $a(n)$: Nó pai do nó $n$.
    \item $\mathcal{S}$: Conjunto de submercados.
    \item $\mathcal{I}$: Conjunto de todos os \textit{trades} (contratos) disponíveis no horizonte.
    \item $\mathcal{I}_n \subseteq \mathcal{I}$: Subconjunto de \textit{trades} cuja data de início é estritamente igual ao mês de ocorrência do nó $n$.
    \item $\mathcal{K} = \{1, \dots, D_{\max}\}$: \textit{Pipeline} temporal de duração dos contratos (em meses defasados), onde $D_{\max} = \max_{i \in \mathcal{I}} (d_i)$.
\end{itemize}

\section{Parâmetros}
\begin{itemize}
    \item $\pi_n$: Probabilidade incondicional de atingir o nó $n \in \mathcal{N}$.
    \item $h_n$: Número de horas no mês correspondente ao nó $n$.
    \item $E_{\text{scale}}$: Fator de escala financeira (ex: $10^6$ para conversão em milhões de R\$).
    \item $C_{\text{inicial}}$: Saldo de caixa inicial da empresa no momento zero.
    \item $L^{\text{cred}}$: Limite mínimo de crédito ou saldo em caixa permitido.
    \item $P_{s,n}$: Preço de Liquidação de Diferenças (PLD) no submercado $s$, nó $n$.
    \item $G_{s,n}$: Geração física consolidada no submercado $s$, nó $n$.
    \item $V_{s,n}^{0B}, V_{s,n}^{0S}$: Volume consolidado (MWm) de contratos legados de compra e venda.
    \item $K_{s,n}^{0B}, K_{s,n}^{0S}$: Preço médio ponderado consolidado (R\$/MWh) dos contratos legados.
    \item $\bar{q}_i^B, \bar{q}_i^S$: Limites máximos permitidos para compra e venda do \textit{trade} $i$.
    \item $K_i^B, K_i^S$: Preços fixos de compra e venda do \textit{trade} $i$.
    \item $d_i$: Duração em meses do \textit{trade} $i$.
    \item $s_i$: Submercado de liquidação do \textit{trade} $i$.
    \item $\lambda$: Fator de aversão ao risco.
    \item $\alpha$: Nível de confiança do CVaR.
\end{itemize}

\section{Variáveis de Decisão e Estado}
\begin{itemize}
    \item $q_{i,n}^B, q_{i,n}^S \ge 0$: Volume comprado e vendido do \textit{trade} $i$ no nó $n$ (MWm).
    \item $v_{s,k,n} \in \mathbb{R}$: \textit{Pipeline} de volume físico do submercado $s$, na defasagem $k$, nó $n$ (MWm).
    \item $c_{s,k,n} \in \mathbb{R}$: \textit{Pipeline} de exposição financeira do submercado $s$, na defasagem $k$, nó $n$ (R\$/h).
    \item $x_n \in \mathbb{R}$: Saldo de caixa financeiro acumulado no nó $n$.
    \item $\mathrm{VaR} \in \mathbb{R}$: \textit{Value-at-Risk} no horizonte final.
    \item $\eta_n \ge 0$: Variável de desvio de perda no nó folha $n$ para cálculo do CVaR.
\end{itemize}

\section{Restrições do Modelo}

\paragraph{Condições de Contorno e Inicialização}
Para o nó raiz (onde $a(n) = 0$), os estados herdados do passado são definidos como nulos, e o caixa é inicializado com o capital preexistente:
\begin{align}
    x_0 &= \frac{C_{\text{inicial}}}{E_{\text{scale}}} \\
    v_{s,k,0} &= 0 \quad \forall\, s \in \mathcal{S},\ k \in \mathcal{K} \\
    c_{s,k,0} &= 0 \quad \forall\, s \in \mathcal{S},\ k \in \mathcal{K}
\end{align}

\paragraph{Limites de Operação e Não-Antecipatividade}
Os volumes negociados estão limitados aos lotes máximos. Para garantir a causalidade temporal e impedir que o modelo acione decisões futuras antes do tempo, as variáveis de \textit{trades} são estritamente nulas se o nó não pertencer ao mês de início do contrato:
\begin{align}
    0 \le q_{i,n}^B &\le \bar{q}_i^B \quad \forall\, n \in \mathcal{N},\ \forall\, i \in \mathcal{I}_n \\
    0 \le q_{i,n}^S &\le \bar{q}_i^S \quad \forall\, n \in \mathcal{N},\ \forall\, i \in \mathcal{I}_n \\
    q_{i,n}^B = 0,\ q_{i,n}^S &= 0 \quad \forall\, n \in \mathcal{N},\ \forall\, i \notin \mathcal{I}_n
\end{align}

Adicionalmente, como o Equivalente Determinístico opera sob um paradigma onde a incerteza é revelada nos nós da árvore, impõe-se a \textbf{Restrição de Não-Antecipatividade Estrita}. Esta restrição garante que decisões tomadas em um determinado momento do tempo sejam idênticas para todos os cenários possíveis que compartilham o mesmo histórico até aquele momento (mesmo nó pai), alinhando a modelagem ao comportamento \textit{Decision-Hazard}:
\begin{align}
    q_{i,m}^B &= q_{i,n}^B \quad \forall\, m, n \in \mathcal{N} \text{ tais que } a(m) = a(n) \\
    q_{i,m}^S &= q_{i,n}^S \quad \forall\, m, n \in \mathcal{N} \text{ tais que } a(m) = a(n)
\end{align}

\paragraph{Transição do Pipeline de Volumes}
A posição comprada ou vendida transita ao longo dos estágios enquanto a duração do contrato for superior à defasagem $k$.
O estado herda a posição do nó pai acrescida das novas contratações do mês vigente:
\begin{align}
    v_{s,k,n} &= v_{s,k+1,a(n)} + \sum_{\substack{i \in \mathcal{I}_n \\ d_i > k,\ s_i = s}} \left(q_{i,n}^B - q_{i,n}^S\right)
    \quad \forall\, s \in \mathcal{S},\ k \in \{1,\dots,D_{\max}-1\},\ n \in \mathcal{N} \\
    v_{s,D_{\max},n} &= \sum_{\substack{i \in \mathcal{I}_n \\ d_i > D_{\max},\ s_i = s}} \left(q_{i,n}^B - q_{i,n}^S\right)
    \quad \forall\, s \in \mathcal{S},\ n \in \mathcal{N}
\end{align}

\paragraph{Transição do Pipeline Financeiro}
De forma análoga aos volumes, a taxa de exposição de caixa (R\$/h) dos contratos travados herda o histórico dos meses anteriores, indexada por submercado:
\begin{align}
    c_{s,k,n} &= c_{s,k+1,a(n)} + \sum_{\substack{i \in \mathcal{I}_n \\ d_i > k,\ s_i = s}} \left(q_{i,n}^S K_i^S - q_{i,n}^B K_i^B\right)
    \quad \forall\, s \in \mathcal{S},\ k \in \{1,\dots,D_{\max}-1\},\ n \in \mathcal{N} \\
    c_{s,D_{\max},n} &= \sum_{\substack{i \in \mathcal{I}_n \\ d_i > D_{\max},\ s_i = s}} \left(q_{i,n}^S K_i^S - q_{i,n}^B K_i^B\right)
    \quad \forall\, s \in \mathcal{S},\ n \in \mathcal{N}
\end{align}

\paragraph{Exposição Físico-Financeira ao Mercado Spot}
Balanço entre a energia física disponível e a entregue por submercado.
A posição física comprada no passado ($k=1$) é injetada no balanço corrente com sinal positivo, provendo energia para liquidação:
\begin{equation}
    E_{s,n} = G_{s,n} + V_{s,n}^{0B} - V_{s,n}^{0S}
              + \sum_{\substack{i \in \mathcal{I}_n \\ s_i = s}} q_{i,n}^B
              - \sum_{\substack{i \in \mathcal{I}_n \\ s_i = s}} q_{i,n}^S
              + v_{s,1,a(n)}
    \quad \forall\, s \in \mathcal{S},\ n \in \mathcal{N}
\end{equation}

\paragraph{Equação de Transição de Caixa}
A variável de estado central intertemporal. O saldo no nó atual herda o saldo do pai acrescido dos fluxos de resultado do mês corrente.
A conversão da base horária ($h_n$) e da escala financeira ($E_{\text{scale}}$) ocorre apenas na injeção de caixa:
\begin{multline}
    x_n = x_{a(n)} + R_n^{\text{leg}}
          + \left[ \sum_{i \in \mathcal{I}_n} \left(q_{i,n}^S K_i^S - q_{i,n}^B K_i^B\right) + \sum_{s \in \mathcal{S}} c_{s,1,a(n)} \right] \frac{h_n}{E_{\text{scale}}} \\
          + \sum_{s \in \mathcal{S}} E_{s,n}\, P_{s,n} \frac{h_n}{E_{\text{scale}}}
    \quad \forall\, n \in \mathcal{N}
\end{multline}
onde $R_n^{\text{leg}}$ é o resultado financeiro determinístico da carteira legada no mês:
\begin{equation}
    R_n^{\text{leg}} = \sum_{s \in \mathcal{S}} \left(V_{s,n}^{0S} K_{s,n}^{0S} - V_{s,n}^{0B} K_{s,n}^{0B}\right) \frac{h_n}{E_{\text{scale}}}
\end{equation}

\paragraph{Restrições de Crédito e Risco}
O saldo não pode violar a margem mínima de crédito estabelecida corporativamente, e o desvio de CVaR é contabilizado exclusivamente sobre os estados terminais:
\begin{align}
    x_n     &\ge L^{\text{cred}}          \quad \forall\, n \in \mathcal{N} \label{eq:credito} \\
    \eta_n  &\ge \mathrm{VaR} - x_n       \quad \forall\, n \in \mathcal{L}
\end{align}

\section{Função Objetivo}

A diretriz ótima é obtida maximizando uma soma ponderada (regulada pelo fator de aversão ao risco $\lambda$) entre a riqueza final esperada e o risco de cauda (CVaR), avaliada apenas sobre os nós folha:
\begin{equation}
    \max \quad \mathbb{E}[x] + \lambda \cdot \mathrm{CVaR}
\end{equation}
\begin{align}
    \mathbb{E}[x]   &= \sum_{n \in \mathcal{L}} \pi_n\, x_n \\
    \mathrm{CVaR}   &= \mathrm{VaR} - \frac{1}{1 - \alpha} \sum_{n \in \mathcal{L}} \pi_n\, \eta_n
\end{align}

% ==============================================================
\section{Implementação via SDDP e Recurso Completo Relativo}
% ==============================================================

\subsection{Topologia do Grafo de Decisão (Policy Graph)}

A formulação multiestágio é estruturada sob a abstração de um \textit{Policy Graph} (Grafo de Políticas), conforme proposto por Dowson e Kapelevich. Definimos o grafo direcionado $\mathcal{G} = (\mathcal{V}, \mathcal{E})$, onde a dinâmica temporal e o processo estocástico são representados pela topologia da rede.

Para este modelo de comercialização de energia, adota-se uma topologia de Grafo Linear, refletindo a progressão sequencial e determinística dos meses do horizonte de planejamento:
\begin{itemize}
    \item \textbf{Vértices (Estágios):} O conjunto de nós $\mathcal{V} = \{1, 2, \dots, T\}$ representa cada mês do horizonte de simulação.
    \item \textbf{Arestas (Transições):} O conjunto de arestas direcionadas $\mathcal{E} = \{(t, t+1) \mid t = 1, \dots, T-1\}$ representa a passagem irreversível do tempo.
    \item \textbf{Dinâmica de Transição:} A probabilidade de transição entre nós é determinística. Estar no nó $t$ implica, com probabilidade $1$, a transição para o nó $t+1$.
\end{itemize}

Nesta arquitetura, a incerteza do problema não governa a transição entre os nós da rede, mas atua como uma variável aleatória $\omega_t \in \Omega_t$ (composta pelas realizações de Geração e PLD) revelada localmente no nó $t$, sob o paradigma \textit{Decision-Hazard}. O agente decisor observa o estado transferido pelo nó $t-1$, toma a decisão de \textit{trading} e, em seguida, o cenário $\omega_t$ se materializa, definindo o fluxo de caixa local e o estado resultante que será transmitido ao nó $t+1$.

\subsection{Recurso Completo Relativo e a Variável de Folga $\sigma_t$}

\paragraph{Definição formal.}
Um problema de programação estocástica multiestágio satisfaz a propriedade de \textbf{recurso completo relativo} (\textit{relatively complete recourse}) se, para todo estado $x_t$ atingível a partir do estado inicial e para toda realização $\omega_t \in \Omega_t$, o subproblema do estágio $t$ admite solução viável. Formalmente, seja $\mathcal{X}_t$ o conjunto de estados atingíveis no estágio $t$; exige-se que
\begin{equation}
    \forall\, x_t \in \mathcal{X}_t,\quad \forall\, \omega_t \in \Omega_t : \quad \mathcal{F}_t(x_t, \omega_t) \neq \emptyset,
\end{equation}
onde $\mathcal{F}_t(x_t, \omega_t)$ denota o conjunto viável do subproblema condicionado ao estado $x_t$ e à realização $\omega_t$.

\paragraph{Por que o DEQ não requer modificação.}
No Equivalente Determinístico, a restrição de crédito~\eqref{eq:credito} participa de um único programa linear de grande porte que engloba todos os nós da árvore simultaneamente. O solver pode, portanto, ajustar as decisões de \textit{trading} em estágios anteriores de modo a prevenir qualquer violação futura. A variável $x_n$ é declarada livre ($x_n \in \mathbb{R}$) e a restrição $x_n \ge L^{\text{cred}}$ é uma desigualdade global do LP, não um limite rígido sobre uma variável de estado isolada. A viabilidade do DEQ é, assim, uma propriedade emergente da otimização conjunta sobre toda a árvore.

\paragraph{Por que o SDDP requer a folga.}
No algoritmo SDDP, cada subproblema do estágio $t$ é resolvido de forma independente, com o estado de entrada $x_t$ fixado pelo passo \textit{forward}. A equação de transição de caixa
\begin{equation}
    x_{t+1} = x_t + R_t^{\text{leg}} + \left[\,\cdot\,\right]\frac{h_t}{E_{\text{scale}}} + \sum_{s \in \mathcal{S}} E_{s,t}\,P_{s,t}\,\frac{h_t}{E_{\text{scale}}}
\end{equation}
contém o termo estocástico $P_{s,t}(\omega_t)$ e $G_{s,t}(\omega_t)$, que não são controláveis pelas decisões $q^B, q^S$ de forma irrestrita (ambas são limitadas superiormente). Para uma realização adversa $\omega_t$, pode ocorrer que nenhuma escolha admissível de $(q^B, q^S)$ satisfaça simultaneamente a igualdade de transição e o limite inferior do estado $x_{t+1} \ge L^{\text{cred}}$, tornando o subproblema inviável e interrompendo o algoritmo.

\paragraph{Formulação com folga (implementação SDDP).}
Para garantir recurso completo relativo, introduz-se a variável de folga $\sigma_t \ge 0$, interpretada economicamente como uma linha de crédito emergencial acionada apenas quando os fluxos de caixa realizados são insuficientes para manter o saldo acima do limite corporativo. A equação de transição de caixa passa a ser:
\begin{equation}
    x_{t+1} = x_t + R_t^{\text{leg}}
              + \left[ \sum_{i \in \mathcal{I}_t} \left(q_{i,t}^S K_i^S - q_{i,t}^B K_i^B\right) + \sum_{s \in \mathcal{S}} c_{s,1,t} \right] \frac{h_t}{E_{\text{scale}}}
              + \sum_{s \in \mathcal{S}} E_{s,t}\, P_{s,t} \frac{h_t}{E_{\text{scale}}}
              + \sigma_t
    \label{eq:transicao_sddp}
\end{equation}
com $\sigma_t \ge 0$ e sem limite inferior rígido sobre o estado $x_{t+1}$. A função objetivo do estágio $t$ é penalizada pelo uso da folga com coeficiente $M \gg 1$:
\begin{equation}
    f_t =
    \begin{cases}
        x_{T+1} - M\,\sigma_T & \text{se } t = T \\
        -M\,\sigma_t           & \text{se } t < T
    \end{cases}
    \label{eq:obj_sddp}
\end{equation}

\paragraph{Consistência entre DEQ e SDDP.}
Seja $x^*$ a solução ótima do DEQ, com $x_n^* \ge L^{\text{cred}}$ para todo $n \in \mathcal{N}$. Nessa solução, $\sigma_t = 0$ para todo $t$, pois a restrição de crédito nunca é violada. Portanto, a política ótima do SDDP com penalidade $M$ suficientemente grande coincide com a política ótima do DEQ: o valor ótimo convergido pelo SDDP é idêntico ao valor ótimo do DEQ, e a folga $\sigma_t$ permanece nula ao longo de toda trajetória ótima. A adição de $\sigma_t$ é, assim, uma modificação de implementação — necessária para garantir a viabilidade dos subproblemas durante o treinamento — sem alterar o modelo matemático formal nem o valor da solução ótima.

\end{document}
